\documentclass[11pt]{article}

% -----------------------------------------------------------------------------
% Reporte mejorado: Validacion futura de TOI/ATI
% No compilar automaticamente desde el pipeline. Compilar manualmente si se desea.
% -----------------------------------------------------------------------------
\usepackage[utf8]{inputenc}
\usepackage[T1]{fontenc}
\usepackage{lmodern}
\usepackage[spanish,es-nodecimaldot]{babel}
\usepackage[a4paper,margin=2.5cm]{geometry}
\usepackage{amsmath,amssymb,booktabs,graphicx,float,hyperref,array,longtable,enumitem,xcolor}

\hypersetup{colorlinks=true,linkcolor=black,citecolor=black,urlcolor=blue}
\setlist[itemize]{topsep=0.25em,itemsep=0.15em}
\setlist[enumerate]{topsep=0.25em,itemsep=0.15em}

\newcommand{\figroot}{../../outputs/toi_ati_future_validation/figures_pdf}
\newcommand{\tabroot}{../../outputs/toi_ati_future_validation/tables}

% Inclusion robusta: si una figura aun no existe, el documento compila con una caja.
\newcommand{\safeincludegraphics}[2][]{%
  \IfFileExists{#2}{%
    \includegraphics[#1]{#2}%
  }{%
    \fbox{\begin{minipage}{0.86\textwidth}\centering
    \vspace{1em}
    Figura no encontrada en la ruta esperada:\\[0.4em]
    \texttt{\detokenize{#2}}\\[0.4em]
    \small Copiar los PDF de figuras a \texttt{outputs/toi\_ati\_future\_validation/figures\_pdf/} o ajustar \texttt{\string\figroot}.
    \vspace{1em}
    \end{minipage}}%
  }%
}

\title{Validación futura de TOI/ATI: estabilidad, ranking conservador y prioridad observacional}
\author{Proyecto Mapper/TDA de exoplanetas}
\date{}

\begin{document}
\maketitle

\begin{abstract}
Este reporte interpreta la fase de validación futura de los índices TOI/ATI. El objetivo no es proponer otro índice principal ni afirmar descubrimientos de exoplanetas, sino auditar la estabilidad del ranking ya construido. La contribución de esta fase es separar anclas con déficit local estable de anclas sensibles al radio, construir una lectura conservadora de ATI y transformar el ranking exploratorio en una priorización observacional más prudente. El resultado principal es que \texttt{HIP 97166 c / cube12\_cluster0} era muy fuerte bajo ATI original, pero pierde fuerza al penalizar sensibilidad al radio; en cambio, \texttt{HIP 90988 b / cube17\_cluster2} emerge como caso más defendible bajo ATI conservador y como primera prioridad observacional. La conclusión metodológica es que TOI/ATI debe leerse como sistema de priorización topológica, no como detección de planetas ausentes.
\end{abstract}

\section{Glosario y notación}

\begin{description}[leftmargin=3.4cm,style=nextline]
    \item[Mapper] Grafo construido a partir de cubiertas del espacio de lentes y clusters locales. Cada nodo representa una comunidad local de planetas. Por el solapamiento de la cubierta, un planeta puede aparecer en más de un nodo.
    \item[Nodo Mapper $v$] Comunidad topológica con miembros $S_v$. En este reporte, un nodo se evalúa como región candidata a incompletitud observacional.
    \item[Planeta ancla $p^*$] Planeta representativo dentro de un nodo. Se usa para estudiar déficit local de vecinos compatibles.
    \item[$R^3$ físico-orbital] Espacio definido por $x_i=(\log_{10}M_{p,i},\log_{10}P_i,\log_{10}a_i)$, donde $M_p$ es masa planetaria, $P$ es periodo orbital y $a$ es semieje mayor.
    \item[TOI] \textit{Topological Observational Incompleteness}. Prioridad regional de un nodo Mapper.
    \item[ATI] \textit{Anchor Topological Incompleteness}. Prioridad local de un planeta ancla.
    \item[$\Delta_{\mathrm{rel,best}}$] Máximo déficit relativo positivo entre radios locales. Es útil para exploración, pero puede sobrepriorizar casos sensibles a la escala.
    \item[ATI conservador] Variante interpretativa que penaliza sensibilidad al radio y radios grandes negativos. No reemplaza TOI/ATI original; los audita.
\end{description}

\section{Problema que resuelve esta validación}

La fase anterior ya produjo rankings TOI/ATI. Sin embargo, un ranking alto no implica automáticamente evidencia robusta. En particular, el ATI original puede subir por un valor alto de $\Delta_{\mathrm{rel,best}}$, aunque ese déficit aparezca solo en un radio pequeño y desaparezca, o incluso cambie de signo, cuando se evalúan radios más amplios.

El problema metodológico es:
\[
\Delta_{\mathrm{rel,best}}(p^*)=\max_r \Delta_{\mathrm{rel}}(p^*,r)
\]
resume el caso más favorable para el ancla. Eso es útil para no perder candidatos, pero no basta para afirmar estabilidad topológica.

Por eso esta fase introduce una pregunta más estricta:
\[
\text{¿el déficit local permanece cuando cambia la escala del vecindario?}
\]

La respuesta se evalúa comparando tres radios:
\[
r_{\mathrm{kNN}},\qquad r_{\mathrm{node\_median}},\qquad r_{\mathrm{node\_q75}}.
\]

Si el déficit es positivo en los tres radios, el caso es más estable. Si solo aparece en $r_{\mathrm{kNN}}$ y se vuelve negativo en radios mayores, el caso es sensible a la escala y debe presentarse como exploratorio.

\section{Marco metodológico}

\subsection{Índice regional TOI}

El índice regional se mantiene como:
\[
\mathrm{TOI}(v)=\mathrm{Shadow}(v)(1-I_{R^3}(v))C_{\mathrm{phys}}(v)S_{\mathrm{net}}(v).
\]

La interpretación de cada factor es:
\begin{itemize}
    \item $\mathrm{Shadow}(v)$ mide frontera observacional local.
    \item $1-I_{R^3}(v)$ penaliza dependencia de imputación en masa, periodo y semieje mayor.
    \item $C_{\mathrm{phys}}(v)$ favorece continuidad física con los vecinos.
    \item $S_{\mathrm{net}}(v)$ favorece soporte de red y penaliza nodos aislados o pequeños.
\end{itemize}

Un TOI alto significa que el nodo es una región prioritaria para inspección. No significa que la región contenga planetas no observados confirmados.

\subsection{Índice de ancla ATI original}

El ATI original se define como:
\[
\mathrm{ATI}(p^*)=\mathrm{TOI}(v)\Delta_{\mathrm{rel,best}}(p^*)(1-I_{R^3}(p^*))A(p^*),\qquad p^*\in S_v.
\]

Este índice combina prioridad regional, déficit local máximo, baja imputación del ancla y representatividad del planeta dentro del nodo. Su punto débil es que $\Delta_{\mathrm{rel,best}}$ puede ser demasiado optimista.

\subsection{Déficit relativo por radio}

Para un ancla $p^*$ y radio $r$:
\[
B_r(p^*)=\{q:\|x(q)-x(p^*)\|\leq r\}.
\]

El número observado de vecinos compatibles es:
\[
N_{\mathrm{obs}}(p^*,r)=|B_r(p^*)|,
\]
excluyendo al propio ancla. La referencia local esperada es $N_{\mathrm{exp}}(p^*,r)$. Entonces:
\[
\Delta N(p^*,r)=N_{\mathrm{exp}}(p^*,r)-N_{\mathrm{obs}}(p^*,r),
\]
\[
\Delta_{\mathrm{rel}}(p^*,r)=\frac{N_{\mathrm{exp}}(p^*,r)-N_{\mathrm{obs}}(p^*,r)}{N_{\mathrm{exp}}(p^*,r)+\epsilon}.
\]

Valores positivos indican menos vecinos observados que esperados bajo la referencia local. Valores negativos indican más vecinos observados que esperados. Ninguno de los dos debe interpretarse como número real de planetas ausentes; son referencias topológicas locales.

\subsection{ATI conservador}

La validación construye variantes interpretativas:
\begin{align*}
\mathrm{ATI}_{\mathrm{stable}} &\approx \mathrm{ATI}\times \text{factor de estabilidad por radio},\\
\mathrm{ATI}_{\mathrm{radius\_penalized}} &\approx \mathrm{ATI}\times \frac{\max(0,\overline{\Delta}_{\mathrm{rel}})}{\Delta_{\mathrm{rel,best}}+\epsilon},\\
\mathrm{ATI}_{\mathrm{conservative}} &\approx \mathrm{ATI}_{\mathrm{stable}}\times \text{penalización por radios grandes negativos}.
\end{align*}

La idea no es reemplazar el ATI original, sino separar dos usos:
\begin{itemize}
    \item ATI original: útil para descubrir candidatos exploratorios.
    \item ATI conservador: útil para priorizar casos defendibles en exposición y trabajo futuro.
\end{itemize}

\section{Resultados: casos estables vs sensibles}

\begin{figure}[H]
\centering
\safeincludegraphics[width=.82\textwidth]{\figroot/stable_vs_sensitive_deficit.pdf}
\caption{Casos estables vs sensibles al radio. El eje horizontal resume el déficit relativo medio y el eje vertical mide sensibilidad al cambio de radio. Los casos con sensibilidad alta dependen mucho de la escala; los casos cercanos a sensibilidad baja son más robustos.}
\label{fig:stable-sensitive}
\end{figure}

\paragraph{Interpretación de la Figura~\ref{fig:stable-sensitive}.}
La figura separa candidatos según dos dimensiones: magnitud promedio del déficit y variabilidad entre radios. Esta separación es crucial porque un caso puede tener un déficit local favorable en una escala, pero volverse sobrepoblado en otra. Los puntos con déficit medio negativo no deben considerarse candidatos de submuestreo local; representan regiones donde hay más vecinos observados que esperados bajo la referencia usada. Los puntos con sensibilidad alta son frágiles: dependen de la definición de radio. Los puntos con déficit positivo y baja sensibilidad son los mejores candidatos para una lectura conservadora.

\paragraph{Observación técnica.}
La escala de la figura puede quedar dominada por outliers con déficit medio muy negativo o sensibilidad extrema. Para una versión final de presentación convendría producir una segunda figura con límites acotados alrededor de cero, por ejemplo $-0.5\leq\overline{\Delta}_{\mathrm{rel}}\leq0.5$, para no comprimir los casos interesantes.

\begin{figure}[H]
\centering
\safeincludegraphics[width=.82\textwidth]{\figroot/deficit_profiles_selected_anchors.pdf}
\caption{Perfiles de déficit relativo por radio para anclas seleccionadas. La estabilidad se evalúa observando si $\Delta_{\mathrm{rel}}$ permanece positivo en $r_{\mathrm{kNN}}$, $r_{\mathrm{node\_median}}$ y $r_{\mathrm{node\_q75}}$.}
\label{fig:deficit-profiles}
\end{figure}

\paragraph{Interpretación de la Figura~\ref{fig:deficit-profiles}.}
Esta es una de las figuras más informativas del reporte. \texttt{HIP 97166 c / cube12\_cluster0} muestra un patrón sensible al radio: el déficit es positivo en $r_{\mathrm{kNN}}$, pero se vuelve negativo en radios más amplios. Esto significa que su prioridad original dependía de una escala local estrecha. Por tanto, es un buen ejemplo de ancla exploratoria, pero no el caso más estable.

En contraste, \texttt{HIP 90988 b / cube17\_cluster2} mantiene déficit positivo en los tres radios. Su déficit no es grande, pero es consistente. Por eso es más defendible bajo una lectura conservadora. \texttt{HD 4313 b / cube17\_cluster10} también parece importante porque presenta déficit positivo en todos los radios y una magnitud relativa mayor en el radio mediano. \texttt{HD 42012 b} aparece en dos nodos, lo que lo vuelve útil para explicar solapamiento Mapper y transición topológica.

\paragraph{Conclusión de esta sección.}
El resultado no dice que \texttt{HIP 90988 b} o \texttt{HD 4313 b} tengan planetas faltantes confirmados. Dice que, bajo la referencia topológica local, son anclas más estables que \texttt{HIP 97166 c} para estudiar submuestreo físico-orbital.

\section{Resultados: cambio de ranking}

\begin{figure}[H]
\centering
\safeincludegraphics[width=.82\textwidth]{\figroot/ati_vs_ati_conservative.pdf}
\caption{ATI original frente a ATI conservador. La línea diagonal indica igualdad aproximada entre ambos. Puntos por debajo de la diagonal pierden prioridad al penalizar sensibilidad al radio.}
\label{fig:ati-vs-conservative}
\end{figure}

\paragraph{Interpretación de la Figura~\ref{fig:ati-vs-conservative}.}
La figura muestra el efecto central de la validación. \texttt{HIP 97166 c / cube12\_cluster0} aparece con ATI original alto, pero cae fuertemente en ATI conservador. Eso confirma que el índice original lo priorizaba por una señal local favorable que no se conserva al ampliar el radio. En cambio, \texttt{HIP 90988 b / cube17\_cluster2} queda alto en la versión conservadora, lo que sugiere una prioridad más defendible.

La lectura metodológica es:
\[
\boxed{\text{ATI conservador corrige sobrepriorizaciones generadas por }\Delta_{\mathrm{rel,best}}.}
\]

\begin{figure}[H]
\centering
\safeincludegraphics[width=.82\textwidth]{\figroot/rank_shift_after_stability_penalty.pdf}
\caption{Cambio de ranking tras penalizar sensibilidad al radio. La comparación entre ranking original y ranking conservador identifica casos que bajan por inestabilidad y casos que se mantienen por robustez.}
\label{fig:rank-shift}
\end{figure}

\paragraph{Interpretación de la Figura~\ref{fig:rank-shift}.}
El cambio de ranking debe leerse como auditoría, no como contradicción. Los casos que bajan no son necesariamente falsos; son casos exploratorios cuyo argumento dependía de una escala local. Los casos que suben o se mantienen no necesariamente tienen déficit grande; tienen déficit más estable o una combinación más defendible de factores. Esta figura es útil para separar tres tipos de anclas:
\begin{enumerate}
    \item \textbf{anclas exploratorias}: ATI original alto, ATI conservador bajo;
    \item \textbf{anclas robustas}: ATI original y conservador altos;
    \item \textbf{anclas de transición}: valor interpretativo por membresía múltiple, aunque su déficit sea pequeño.
\end{enumerate}

\section{Casos finales para trabajo futuro}

\begin{figure}[H]
\centering
\safeincludegraphics[width=.9\textwidth]{\figroot/final_future_work_cases.pdf}
\caption{Casos finales para trabajo futuro. La figura resume cinco funciones interpretativas: región top por TOI, ancla top por ATI original, ancla top por ATI conservador, ancla repetida y ancla con déficit estable.}
\label{fig:future-cases}
\end{figure}

\paragraph{Interpretación de la Figura~\ref{fig:future-cases}.}
La selección de casos finales es útil porque convierte el ranking en una narrativa. Sin embargo, hay dos puntos que deben corregirse o verificarse antes de la versión final.

Primero, el caso \texttt{top\_toi\_region} aparece como \texttt{nan} en la figura. Conceptualmente debería corresponder a \texttt{cube12\_cluster0}, que es la región con mayor TOI. Esto parece un problema de unión de tablas: una región no tiene \texttt{anchor\_pl\_name}, por lo que el campo de visualización queda vacío. La solución es usar campos separados para regiones y anclas, por ejemplo \texttt{display\_name}, \texttt{node\_id} y \texttt{anchor\_pl\_name}.

Segundo, el caso \texttt{stable\_deficit\_anchor} aparece como \texttt{K2-147 b} con valores nulos de TOI y ATI conservador. Aunque pueda ser estable bajo déficit, no es necesariamente un buen caso de exposición si su prioridad global es cero. Para exposición conviene elegir un caso estable con prioridad no trivial, como \texttt{HIP 90988 b} o \texttt{HD 4313 b}, dependiendo de la tabla final.

\paragraph{Narrativa recomendada de casos.}
Para una exposición, la secuencia más clara sería:
\begin{enumerate}
    \item \textbf{Región top por TOI}: \texttt{cube12\_cluster0}. Explica el índice regional.
    \item \textbf{Ancla sensible al radio}: \texttt{HIP 97166 c}. Explica por qué ATI original puede inflar un caso.
    \item \textbf{Ancla conservadora}: \texttt{HIP 90988 b}. Explica estabilidad bajo radios.
    \item \textbf{Ancla de transición}: \texttt{HD 42012 b}. Explica solapamiento Mapper.
    \item \textbf{Candidato estable alternativo}: \texttt{HD 4313 b}. Explica déficit positivo en radios múltiples.
\end{enumerate}

\section{Prioridad observacional}

\begin{figure}[H]
\centering
\safeincludegraphics[width=.82\textwidth]{\figroot/observational_priority_ranking.pdf}
\caption{Ranking de prioridad observacional futura. El puntaje no implica confirmación de objetos ausentes; ordena anclas donde futuras observaciones o análisis de completitud podrían ser más informativos.}
\label{fig:observational-priority}
\end{figure}

\paragraph{Interpretación de la Figura~\ref{fig:observational-priority}.}
El ranking de prioridad observacional resume la lectura conservadora. \texttt{HIP 90988 b / cube17\_cluster2} aparece como candidato principal, seguido por casos como \texttt{HD 42012 b}, \texttt{HD 11506 d}, \texttt{24 Sex b}, \texttt{HD 4313 b}, \texttt{HD 167677 b}, \texttt{HD 13931 b} y \texttt{Kepler-48 e}. La presencia repetida de algunos planetas en diferentes nodos no debe interpretarse como duplicación errónea: es consecuencia del solapamiento Mapper y puede indicar posición en zonas de transición topológica.

La utilidad de este ranking es práctica. Permite decir: si se buscara priorizar inspección futura, no empezaríamos por el ancla con mayor ATI original, sino por anclas cuyo déficit sea más estable, cuya imputación sea baja y cuya posición topológica sea interpretable.

\section{Discusión general}

Esta fase cambia la interpretación del proyecto. Antes, el argumento era que algunas regiones tenían alto TOI/ATI. Ahora, el argumento es más refinado:
\[
\boxed{\text{TOI identifica regiones; ATI original explora anclas; ATI conservador prioriza casos defendibles.}}
\]

Esto permite separar casos que parecían fuertes por una escala local favorable de casos que conservan señal al cambiar el radio. La consecuencia científica es importante: la evidencia de incompletitud observacional no debe basarse solo en el máximo déficit local, sino en la estabilidad del déficit, la baja imputación, la continuidad física y el soporte de red.

\subsection{Qué cambia frente al reporte anterior}

El reporte anterior ya separaba déficit relativo y absoluto. Esta nueva fase agrega una capa adicional: no basta con calcular $\Delta_{\mathrm{rel}}$ correctamente; también hay que preguntar si el signo y la magnitud de $\Delta_{\mathrm{rel}}$ son estables entre radios.

El caso \texttt{HIP 97166 c} ilustra el problema: es fuerte bajo ATI original, pero sensible al radio. El caso \texttt{HIP 90988 b} ilustra la solución: no tiene el déficit más grande, pero sí una señal más estable. Esto convierte la discusión en una comparación entre magnitud y robustez.

\subsection{Qué significa físicamente}

Si las anclas principales pertenecen a comunidades dominadas por velocidad radial, la dirección esperada de incompletitud no es arbitraria. La hipótesis física prudente es que el catálogo puede estar submuestreado hacia:
\[
\text{menor masa planetaria}
\]
o
\[
\text{menor señal dinámica detectable por velocidad radial}
\]
a periodo y semieje comparable. Esta frase debe mantenerse como hipótesis de priorización, no como predicción de planetas reales.

\subsection{Qué significa topológicamente}

Topológicamente, un caso conservador indica que el planeta ancla vive en una región donde la red Mapper, la geometría en $R^3$ y el déficit por radio apuntan en una dirección coherente. Un caso repetido, como \texttt{HD 42012 b}, puede señalar transición entre vecindarios Mapper. Esto no es un error: es una propiedad esperada de Mapper cuando la cubierta tiene solapamiento.

\section{Conclusiones por nivel de fuerza}

\subsection{Conclusiones seguras}

\begin{enumerate}
    \item La validación por radio mejora la interpretación de ATI.
    \item ATI conservador reduce la influencia de $\Delta_{\mathrm{rel,best}}$.
    \item \texttt{HIP 97166 c} es un caso exploratorio sensible al radio.
    \item \texttt{HIP 90988 b} es más defendible bajo una lectura conservadora.
    \item \texttt{cube12\_cluster0} sigue siendo la región regionalmente más importante por TOI, aunque debe corregirse el campo \texttt{nan} en la figura de casos finales.
\end{enumerate}

\subsection{Conclusiones prudentes}

\begin{enumerate}
    \item Algunas regiones Mapper son candidatas razonables a incompletitud observacional topológica.
    \item Algunas anclas pueden funcionar como puntos de inspección para vecindarios físico-orbitales submuestreados.
    \item La prioridad observacional debe basarse más en ATI conservador que en ATI original cuando el objetivo es exposición científica o seguimiento futuro.
\end{enumerate}

\subsection{Conclusiones que no se deben hacer}

\begin{enumerate}
    \item No se puede afirmar que se descubrieron planetas nuevos.
    \item No se puede afirmar que faltan exactamente $N$ planetas.
    \item No se puede convertir un déficit topológico en una función de completitud instrumental.
    \item No se debe interpretar un caso sensible al radio como evidencia fuerte de submuestreo.
\end{enumerate}

\section{Limitaciones}

\begin{enumerate}
    \item \textbf{Sin función de completitud instrumental.} El análisis no modela $P(\mathrm{detectado}\mid M_p,P,a,\mathrm{estrella},m)$. Por tanto, los déficits son topológicos, no tasas físicas de detección.
    \item \textbf{$R^3$ simplifica la física.} Masa, periodo y semieje mayor no capturan inclinación, excentricidad, masa estelar, ruido instrumental ni cadencia de observación.
    \item \textbf{El proxy RV no es amplitud RV real.} Si se usa un proxy dinámico, debe describirse como heurístico.
    \item \textbf{Mapper permite membresía múltiple.} Esto es metodológicamente normal, pero debe explicarse para no confundir planetas repetidos con duplicados.
    \item \textbf{El caso final regional contiene un bug de visualización.} El campo \texttt{top\_toi\_region} no debe aparecer como \texttt{nan}; debe mapearse a \texttt{cube12\_cluster0}.
    \item \textbf{Un déficit estable puede ser pequeño.} Estabilidad no equivale a magnitud alta. \texttt{HIP 90988 b} es más estable que \texttt{HIP 97166 c}, pero su déficit sigue siendo débil.
\end{enumerate}

\section{Trabajo futuro}

\subsection{Corregir la tabla de casos finales}

La mejora inmediata es técnica: corregir la fila \texttt{top\_toi\_region} para que muestre \texttt{cube12\_cluster0} en lugar de \texttt{nan}. También debe revisarse la elección de \texttt{stable\_deficit\_anchor} para evitar seleccionar casos con TOI y ATI conservador nulos.

\subsection{Crear fichas interpretativas por caso}

El siguiente reporte debería incluir fichas para tres a cinco casos. Cada ficha debe tener:
\begin{itemize}
    \item nodo Mapper;
    \item planeta ancla, si aplica;
    \item TOI;
    \item ATI original;
    \item ATI conservador;
    \item clase de estabilidad;
    \item déficit por los tres radios;
    \item imputación en $R^3$;
    \item método dominante;
    \item dirección física esperada de incompletitud;
    \item cautela principal.
\end{itemize}

\subsection{Incorporar funciones de completitud instrumental}

La validación más importante sería modelar:
\[
P(\mathrm{detectado}\mid M_p,P,a,\mathrm{estrella},m).
\]
Con esta función, el déficit topológico podría compararse con detectabilidad real por método. Esto permitiría distinguir mejor entre incompletitud física plausible y artefactos de muestreo.

\subsection{Validar con catálogos futuros}

Guardar los rankings actuales y volver a ejecutar el pipeline con versiones futuras del catálogo. La pregunta sería:
\[
\text{¿las regiones con alto TOI/ATI conservador reciben nuevos planetas o mediciones más completas?}
\]
Si esto ocurre, TOI/ATI ganaría valor predictivo como herramienta de priorización.

\subsection{Simulaciones de inyección-recuperación}

Insertar planetas sintéticos en $R^3$, aplicar funciones de selección simuladas por método y evaluar si TOI/ATI recupera las regiones artificialmente incompletas. Esta sería la prueba controlada más fuerte del método.

\subsection{Propiedades estelares y proxy dinámico}

Agregar masa, radio, temperatura efectiva, magnitud y distancia estelar. Para velocidad radial, una versión más física del proxy debería aproximar:
\[
K\propto M_pP^{-1/3}M_{\star}^{-2/3}.
\]
Esto no reemplaza una simulación instrumental, pero mejora la interpretación física.

\subsection{Rankings separados por método de descubrimiento}

Construir rankings separados para velocidad radial, tránsito, imagen directa y microlensing. Las direcciones de incompletitud esperadas son distintas para cada método; mezclarlas en un único ranking puede ocultar mecanismos diferentes.

\subsection{Validar estabilidad Mapper}

Repetir los rankings bajo variaciones de parámetros Mapper: número de cubos, overlap, lente, métrica y clustering. Un candidato que sobrevive a estas variaciones es más fuerte que uno dependiente de una configuración específica.

\section{Conclusión}

La validación futura convierte TOI/ATI de un ranking exploratorio en un sistema de priorización auditado. La principal mejora es distinguir anclas sensibles al radio de anclas más estables. El resultado más importante no es que existan planetas faltantes, sino que el catálogo puede evaluarse topológicamente para decidir dónde buscar evidencia de incompletitud observacional.

La frase final recomendada es:
\[
\boxed{\text{TOI/ATI no detecta planetas ausentes; prioriza regiones donde la incompletitud observacional es topológicamente plausible y auditable.}}
\]

\end{document}
